\subsection*{Source Metadata}
\noindent\textbf{Source File:} \path{How News Becomes News- Workflow & Editing Explained - MCM 202 - Lecture 2.txt}\par
\noindent\textbf{Document Type:} Lecture Transcript\par
\noindent\textbf{Display Title:} MCM 202 - Lecture 2: How News Becomes News - Workflow and Editing Explained\par
\noindent\textbf{Source URL:} \url{https://www.youtube.com/watch?v=lhJDKaUJ5fw}\par

\subsection*{Transcript Text}
{\small
\begin{Verbatim}
Welcome to lesson number two. The newsroom workflow, sub editors, and the editing symbols. Okay. So,
today we'll be looking at some critical aspect of things, and this is probably one of the most
detailed sessions we'll be having. So, I'm still your lecturer, teacher, whatever you call it. If
you want to stay connected, please use the QR code to stay connected. Before we dive into today's
session, you've been with me for a while now, and then I would need you to, you know, engage. Check
in. Tell me, what do you know? What do you know? Check in. Hey, student, this is the way to mark
your attendance to show that you're here. From the last session, how was it? Put it in the comments.

What did you learn from the last session? And check yourself in. All right. Thank you. So, today
we'll be covering something somewhat different, building up to what we did the last time. And we'll
be looking at the newsroom workflow. We'll be looking at the roles and responsibilities of a sub
editor, as well as I'll be introducing you to the basic editing symbols. So, write with me and enjoy
today's session. Okay. Now, first, I would also include some very practical aspect of things, so be
ready for that aspect of things. Some practicals will be playing out in today's teaching. Okay. So,
that being said, let's take a look at the objective of today's session. Today would be explaining
how content move through the newsroom.

We'll be identifying the roles of a sub editor, and then we'll be, you know, applying some basic
editing and providing symbols to our works digitally, particularly digitally, because that's the age
we're living in. And then lastly, we'll be under We'll learn how to understand how to edit how
editing decisions shape the final outputs in practicals. Okay. Once again, if you haven't checked
in, please do so now. Okay. What we're covering today, let's take a look at them once again. What
we're covering today. I've shared it with you. Do well to take, you know, all the things we're
covering today.

And tell me, what do you already know about these things? What do you already know about these
things? Put it in the comments. Put it in the comments. Let me know what you know. Let me know what
you know. Okay. Thank you. Now, I'm just going to take you straight on today's lecture. All right.
Here's my face. I would like to ask you, think about this. Imagine that there is a breaking news
happening right now, right where you are. How does that information go from a reporter's notebook to
your phone screen? What processes does it pass through for it to before it gets to your to your
mobile phone? What are the steps? What do you think? What procedures does it pass through before it
gets to your mobile phone? Does it get there by magic? Or does it pass through some standard of
processes before it gets to your phone? This is our concern today.

And we want to ensure that we'll make the most of it to understand how it works. You see, the news
process is a structured process, and it's required it requires structured approach. So, we'll say
the newsroom workflow is a structured process through which news contents is created, reviewed, and
also published before it gets to your phone. Someone have to vet it. Someone have to ensure that
what you'll be getting is original, is correct, is intact. And it has to be accurate. It should be
clear, at least to the reader. And it should be timeless. no, not really timeless. It has to come to
you in time, that's what I mean. Timeliness, yeah. All right.

So, the basic flow is this. It goes from the reporter to the sub editor, from the sub editor to the
editor, from the editor to what? Publication. Does that make sense? From the reporter to the sub
editor, from the sub editor to editor, and then to the publication. So, before it gets to your
phone, it needs to, at least before after the reporter, two other people must should look through it
before it gets to you. Let's take a look at the stages. Stage one. News gathering. This is the stage
of, you know, gathering, getting the news, and also reporting in the journalist collecting
information. They source, they include interviews, they attend events, they analyze documents. They
go through challenges such that such as answering questions about incomplete informations.

They have to do this within a certain time frame, and ensure that they filter out biases. All right.
Now, the second stage will be drafting. What do they do? They write the initial version, and then
they focus on the facts to ensure that the facts are written correctly. And then they attend to the
common issues such as grammar errors, poor structures, and and they talk about missing contexts. You
see, stage one has to do with, you know, sourcing it. The stage two, writing it. And they have to
ensure, as a reporter, as a trained journalist who who you are becoming or media person who you
become, you need to be able to pick out grammatical errors. Ensure that your structures are good,
and ensure that you're not missing out contents.

Let me test you. Let's see if you've been following. What is the correct order of news workflow? Is
it A, editor, reporter, sub editor, publication? Reporter, editor, sub editor, publication?
Reporter, sub editor, editor, publication? Or sub editor, reporter, editor, publication? The answer.
Could it be this? That is wrong. It can't be that. Think again. What would it be? It goes from the
reporter to who? From the reporter, it goes to somebody. Who is that person? Think about it. Think
about it. From the reporter, who does it go to next? After the reporter. Where in the chat? Where in
the chat? I'm waiting. All right. Thank you.

It goes to the sub editor. The sub editor takes the sub editing process, and then goes to the
editor, and the editor publishes it. You know, sends it to the publication. Now you know this. Okay.
Let's move on. Let's try another one. Which stage is from really responsible for refining and
correcting the news content? Reporter, sub-editing, publishing, broadcasting. For refining and
correcting the news content, after that should literally come after the reporter. Correct,
sub-editing is correct. So, all right. Why is the news flow as important Okay, this is a question.
Why is the newsroom workflow important? The newsroom workflow, why is it as important as even the
news process itself? Is it because it increases the number of stories published? It speeds up
printing only.

It eliminates need for editors, or it ensures accuracy, clarity, and quality control. Okay, the
answer is it eliminates, it ensures accuracy, clarity, and also quality control. Let's take a look
at stage three of this, which is sub-editing. Remember, the reporter gather writes, then sub-edits.
Meaning, in the first step flow, we have the reporter, sub-edit, edit, and then publication. So, the
first major control stage is the sub-editing stage. This is followed by the editing stage. So, the
sub-editor reviews and they refine content. They review and they refine content. And their key
responsibilities include to correct grammar, to you know, correct errors, to ensure that the content
is clear, and they check facts.

They check facts. And the last thing they do is to write headlines and captions. They write
headlines and captions. This takes us to the stage four, which is editorial review. The editorial
review is done by the editor himself and then it's like a final editorial control stage. Okay, what
happens there is that the senior editor reviews the editor and what has been you know, sub-edited.
You understand? They review the content. They look out to ensure that there's alignment with the
editorial policies, ensure that the content is right, the content is appropriate, and their focus
basically is to reduce legal risk, so they don't get into trouble, and also to what? To reduce and
the the chances of breaching ethical codes.

So, ethical consideration is very important at this stage. And of course, to ensure that the news
carries what they call news value. And then this takes us to the next stage of stage five, which is
the publishing stage. The publishing stage, very important. It's this is where the content is sort
of released to you. You know, when I asked what are the stages, what are the things they pass
through before it gets into your phone? You know, these are the stages and then the point of
publishing is where it gets to your phone, where you get notified about the process. Okay. So, the
content is released to the audience. The platforms could be prints, could be magazines, could be
broadcast, it could be digital, like I've said, could be online like that, and like and so on and so
forth. So, what depending on what is being edited. And the key reality is that errors at this stage
affects public trust.

So, you need to ensure that you get your editing right early before it gets to this stage. this is
editing and graphics. And we're sort of looking at the fundamental aspect of things before we get to
the very practical side of things. So, before a graphic design becomes a graphic design, there's
what we call a copy. Before it is designed to become the newspaper layout, there there are processes
that it has gone through, sub-editing and more. So, if an an article is poorly edited, and it makes
its way to the graphics designer, to the layout person, that's is where they'll plan. And you want
to avoid that at all cost. I have a parallel course that that's looks closely to this and it's
currently on YouTube on that courses and then it's on layout and design.

They are sort of gliding in scope with this, but you need to know that for editing and graphics, you
you can't afford to have an error. Whatsoever error you have is what will make it into the graphics.
What is editing and graphics? We know editing is editing, like I shared in the first first lesson.
Addition and and subtraction. And graphics is of course the use of graphs, image, pictures,
illustrations to communicate a message. So, when we have a copy, a copy is like a brief that would
drive that way from the design. So, we're going to get there, but for now, what is the primary
responsibility of a sub-editor? Is it gathering news from sources, publishing content online,
reviewing and checking content before publication? Printing newspapers, which of this? The answer is
C, reviewing, refining, and checking content for publication.

All right, next question. Which of the following should a sub-editor prioritize? What should it be?
Verifying facts and accuracy, correcting punctuations, writing catchy headlines, or adjusting font
styles? Which is very important? Yeah, you could correct punctuations and still have unverified
facts. You could write a catchy headline and still have unverified facts. No, and then adjust font
size but full of errors. So, the answer will be verifying facts and accuracy. And this is very
important, so you don't get into trouble as a news organization. Okay, let's take a look at this.
Which quality is most essential for a sub-editor? Is it physical strength, drawing skills, public
speaking, or attention to details? physical strength? not really.

you don't need that, but you need it as a human. Drawing skills? Oh, that's not necessary for a
sub-editor. Public speaking? We are not speaking here. But, attention to details, that's quite
important. So, the answer is attention to details. All right, let's define who a sub-editor is. A
sub-editor is a professional who is responsible for refining content before publication. Yeah,
that's what who they are. They they they work between the reporter and of course the editor because
they are like the linkage. So, a sub-editor is that person who you report to immediately after
writing something. And you hand over your content to the sub-editor to sort of refine, review,
improve.

Before it gets to the editor, and the editor makes a final conclusion. So, I wanted to I want us to
do a a checking now. Let's consider this. And we're going to do this for about six to eight minutes.
And I want you to think like a sub-editor. A reporter submits a story with poor grammar to you. And
what would be your first step and why why why would you take that step? He submits a story. It's as
a poor grammar or it has some poor grammar, weak headlines, unverified facts. What would you do
first? And why would you do what you do? This is a task for you. I need to do this now. You know,
you I want you to you know sort of on your Telegram group engage one another. Ask replies, respond
to one another.

It shouldn't be a one-way thing. If you're here on on YouTube, let's engage. Ask one another,
respond to to chat and more. So the question is let's have a discussion around this. What will be
your first action? What will be your second action? What will be your justification? So I'll stay on
this slide for a while. A reporter submits a story to you and you can see that there are poor
grammar in there and then there are weak headlines in there. There are unverified facts in there.
What would you do first and why? Put it in the chat. What would you do and why? Don't disengage.
This is a time for you to think. This is a time for you to write. This is a time for you to reflect.
Don't just be idle. Don't fast forward either. Do something now. I'm sure you engage, please.

Do it. What would you be your first action? What would be your second action? What would be your
justification? Put it in the chat. Thank you for sharing that. now, let's take a look at the core
responsibilities of the sub editors. Now, the core responsibilities of the sub sub editors include
the fact that they edit for grammar and clarity. They fact-check content. They ensure consistency
with style guides. They write headlines and captions. They detect legal and ethical risk in content.

And they also have some skills that make them quite very effective. What are the skills? The skills
are that they pay attention to details. They give an uncommon attention to details. They try as much
as possible to attend to any sort of errors that could the content. So they have strong language
skills, critical thinking, speed and accuracy, ethical awareness. And all of these put together
makes them effective at what they do. Okay. Now, typical errors in news copy which are the sub
editors sort of catch include grammar mistakes repetitions, ambiguity You know what repetition is
when you repeat things. Ambiguity when it's confusing and it's not very clear. Inaccurate facts.

Poor sentence structures. These are some of the things that the sub editor looks at. Moving on. They
editing symbols are also very important. These are shorthand way of correcting without having to
write sentences. If you want to know more about this look out for the video I'll put in the
description. I'll particularly timestamp the the moment where I talked about it in another course so
I don't have to repeat this here. So we have the carrot signs, we have the insert signs, we have the
transpose sign, we have the capitalization sign, we have the the the bolding sign. We have the stet
sign meaning that keep it in place. So these are used in editing, subediting and more. However,
today I'll be describing and focusing more on the digital way of doing things.

So why are these editing symbols important? They help provide clear instructions for corrections.
They save time in manual editing and they ensure consistency across editing across editors sort of
this. So, when Mr. A is using the same sign as B is using the same sign, you're not having troubles
comprehending what they're trying to say. So, it's like a uniform thing. Okay. Attempt this
question. What is the main purpose of editing symbols? To decorate text, to replace editors, to
provide clear corrections, to provide clear corrections instructions, to write headlines. Which of
these? To provide clear corrections instructions. Yeah, that's it. All right.

Now, the key editing symbols like I've mentioned includes delete, insert, transpose, capitalize,
lowercase, new paragraph, you know, bolding and all of these things. These details are also
available using the link in the description. And if you want to read more, there is a chapter in my
monograph that addresses up way in the description, too. Thank you. All right. Let's take care of
this. Transition from digital to Well, let's transit to digital editing. Digital editing uses track
changes, not necessarily symbols, and this is now happening digitally, which I'm going to describe
in a moment for you guys so you to take advantage and know more about.

So, what we do here is that we're going to, you know, learn about the tools in Microsoft Word. You
could use Google Docs for the same purpose, but today I'll be focusing on Microsoft Word. But then
before then, let's take a look at this. Which tool is commonly used today? Instead of manual editing
symbols. Is it typewriter, Microsoft Word track changes, calculator, printer? We You know the answer
already. So, it's Microsoft Word track changes. Don't forget that with Microsoft Word track changes,
which is what we're going to describe in a bit. And then which of the following is an example of an
editing action? Is it publishing a story, intervening, the self-correcting grammatical errors,
printing errors? This sounds like a silly question, but it's correcting grammatical errors. So, all
right.

Let's do something. Now, all right. I'm going to provide you with a flowed paragraph. I'm going to
mark the corrections using editing symbols, and then we'll do something interesting. So, before we
do that, I would like to engage you with some hands-on exercise. did did did Let's move on. To the
hands-on exercise. On the next slide, I need to look at this particular poor written paragraph. And
attempt, make effort to correct it. Do that. Do that. Try. Can you see Can you pick the errors? Can
you see the mistakes? Can you see the grammatical errors? Can you see the agreement errors? Can you
see the capitalization errors? Yeah.

Can see punctuation errors in there, clarity errors, styles errors. Attempt this. I'll stay on here
for about 2 more minutes, and then let's see before we press on. Attempt it. Attempt it. Yes. If you
like, break it in paragraph if that makes sense. Check out for the capitalization errors, the proper
nouns, the commas, the context issues. You can rewrite some things, paraphrase it. Check out for
sharpness, lack of precision. Look Look for wordiness when it is too wordy. Look for repetition. For
example, weak phrases such as yesterday announce instead of announced yesterday.

Look for the spell Look out for the spelling of Nigeria. Look out for for out They said teachers is
That's a subject-verb disagreement. Look out for the is planning. That's a grammatical error. Okay.
Let's bring this home. All right. This is a text here. And I'll be doing the same thing here. And
you know, giving you some very a real-time edits. For you to do this, go to your Microsoft Word,
click on track changes. You toggle it on and off. When it is somewhat gray, it means it's active.
And you can begin editing. Anything you change here will be recorded. So, have a look at it. If you
are here, then you are you are cool. So, the federal federal I can see F.

The F should be in initial caps, not small small case. So, and then the same as the G. Yes. So, the
federal government yesterday announced the Yes. That's they are planning to introduce That's they
are planning to introduce I don't think. I would say do you Federal Government yesterday announced
plans to introduce. You see, when I say they are that they are planning to, it sounds childish. It
sounds childish. Yes. So, introduce a Looking at that, introduce a new education policy that will
improve. That doesn't look good. So, let's see what works here. Yes, the reading back and forth to
see if for the flow, the connection. And this is very important.

So, introduce a new education policy aimed at improving aimed at improving That sounds more like it.
aimed at improving aimed at improving student performance across Nigeria. Full stop. And then the
next paragraph I sort of broke that down. So, it continues from where stopped. Yeah, the education
the Minister education stated stated that See, the one without underline is is going. The one with
underline stays. stated that the policy focuses that is a typographical error. I can't see something
there.

So, it's a very important that one look at things over and over again. Ideally, that should be
stated that the policy focused on or the policies focus on or the policy focuses on digital learning
and teachers training. Full stop. However, many schools lack computer facilities and several
teachers are not properly trained which make or making the making implementation difficult. That
could be another paragraph. Students Capital letter S. in rural area are mostly are the most are
expected to be the most affected. expected to be the most affected. by this these challenges. Okay.
This looks good.

by these challenges. Now, this looks rough, but this is the work of a subeditor. This is what they
do every day. This is what they spend their time and most of their day doing. They do this to
correct errors. They do this to ensure that there is clarity. They do this to ensure that the end
the the meaning is preserved and conveyed. So, have a look at this. I'm still trying to figure out
some things and I'm trying to show you two ways you can delete letter and and then this can affect
the entire work by changing it all. If we highlight a letter and change, it can affect the entire
work. If you just deleted a letter and you insert, it will work for all. So, this is what it looks
like.

And remember I'm using track changes, that's why you're seeing the yellows is blah blah blah. This
how editing is done digitally. Okay. This is another sort of demonstration of how things sort of
work. Mhm. Take a look at this. I'm trying to show you the function track changes, the options that
are available. We have the all markups there. Review review and pan. I'm just clicking on the no,
the all markups and changing to simple. This is the final edit. Even in this final edit, I can still
see an error. This is what our original looks like. And we have the markups which I can change to
show me simple markup. I can change it to give me another option of the entire sort of markup with
the balloon which you can see everything I've changed.

You have the reviewing pan. It can come on the right or at the bottom depending on what I want.
That's the one on the left, sorry. And then we have option for the bottom. just to show things if
you want to to give some more details. So, apart from this you have the one for the left hand side
which sort of is very compact. I think I like that way I need to use it. And then this is where
you're showing markup. Markups are your corrections. And you can of course, you know, accept changes
reject changes. and in this instance, I would like to first save my document. Let's save documents.
So, we have it as a working document. So, you won't see the pop up, but you see the file name
changed.

Doc one is going to change to whatsoever I determine it to be. All right. So, it's saving already.
And it takes a bit of time because it's saving to cloud. Okay. So, that's the original markup. And
then and then with the original text, sorry. And then we have the review that I've shown you
earlier. Then you can come to track changes. You can see the changes that has been done. You can
accept, you can reject. So, we're just going to accept a few changes now. Watch that. Looking at
that F, you see that? I've accepted it. The government, the G in the government, I've accepted it.
So, as you are accepting, it is leaving. It is reducing.

So, and you can as well decide to accept all. See that's another error there. Plans to there's a
space in there. It's not supposed to be So, I've edited it and I have to accept it again. So, now it
needs to move on with that until you do all. But then I want to accept all at a go. So, that means
selecting selecting all. You really can't see that pop up. So, please engage that in your Word
document and see you see the option. Under the accept this in there is accept changes one after the
other, accept and stop tracking. And there is accept all and stop tracking too. I'll accept all. And
yeah, I'm moving on, you know. So, depending So, look at that. Stated that you see why there's a
different levels of editing.

There is subediting, there's main editing. At my subediting, I I probably didn't catch the fact that
I need to put policies at policy. But I think I noted it when I was talking. Just for This is for
teaching purpose only. So, stated that the policies focuses, that is wrong. It should be policies
focus or policy focuses. You see that policy is focuses or policies focus. So, in this case, I'll
stick with policies focus. Yeah, this one is correct if there are more than one policy. If it is
just a single policy, I could say policy focuses. So, now you can see the original. You can see the
no markup. And then I can show you more things around this and what it looks like eventually.

All right. Thank you for learning that with me. Do some practice, too. So, this is what it looks
like. this is what it looks like originally. And yeah, these are the corrections. This side all
markup and when you select, this is what it looks like when the subeditor has done some things or
the editor. And that's is what it looks like below it, the finished work. You understand? Showing it
using simple markup. And then we have the one with no markup versus the original. You see, look at
look at the difference. And look at the works that have been put into this. This is you need to
start learning how to do this as an editor to ensure that you're sort of, you know, doing the right
thing.

All right. So, this is the options I showed you where you can see all revisions. And then on the
bottom and then on the side. And then once again, look at what I've got. Look at what I've got.
That's the edited version up there. and then addition below. And here is the edited version on my
slide. Okay. What have we done here? We change announce, which is an incorrect tense. We changed it
to announced. We change there is planning. It's not This is they are. And then policy of focus. No,
policy focuses. Or policies focus. So, school does not These are some of the things you can read
through that and then mix it away. Thank you. And "Nigeria" should be in capital letter N, not small
letter N, which is what we've done.

All right, let's take a quick quiz. What is the first step in effective editing? Is it writing the
headline? Publishing the content? Identifying and understanding errors? Or adding images? Effective
editing. Will it be all about cutting headline? Will it be all about publishing content? Will it be
identifying and understanding errors? Or will it be adding images? The answer is identifying and
understanding errors. Okay. Let's move on. Why is it important to break long sentences into during
editing? Yeah, it's abs Is it to avoid punctuation? To make text longer? To reduce printing cost? Or
to improve clarity and readability? The answer will be to improve clarity and readability. Okay,
what's the main goal of editing? To increase word count? To impress the editor? To make text
complex? To improve clarity and accuracy and meaning? I believe that is the answer D.

So, today we have been able to learn learn some basic things and newsroom workflow is a short
structured communication. Sub editing are critical to quality control. Editing symbols, standardized
corrections, and then editing improve clarity, improve accuracy, and it also improve credibility. I
have an assignment for you. I Yes, I have an assignment for you. So. I want you to roughly draft a
short news article with errors. I should add the errors There's no way you can ever draft two
without an error. And don't use AI. You can't, of course, use AI to draft a rough or a rough edits,
a rough draft of any story about anything.

Just for you to see. And then edit the the content that has been generated with mistakes. Apply at
least five editing corrections symbols, which I've sort of shared. And then submit both. Use the QR
code, scan it, then use submit. So, you can determine your text yourself. You can use the AI to
generate one wrong thing. And I know some of you got AI to use AI to correct it. Just to ensure that
it is bolded that you are able to pick the errors and know it. My goal is to take you through the
process for you to know. It's your choice to take responsibility or to be responsible. So, learn
responsibly using AI. Okay. Let's construct the news flow. How does it work? How does the new flow,
news flow work? So, I want you to arrange this in the correct order.

We did this earlier. Which is the first one? Which is the second one? Which is the third one? And
which is the last one? All right, beautiful. Thank you for sharing that. It's from the reporter to
the sub editor, sub editor to the editor, editor to publication. We did this earlier. All right.
Now, I want you to sort of put yourself in a position of reflection. Think about it. Engage yourself
in the chat. I've actually provided you the answer. Why is it like that? Why should it be that way?
Why shouldn't be reporter, editor, sub editor after a sub is coming after editor and then
publication? So, I'm going to leave you for another sort of 2 minutes.

Ask yourself. Arrange it in the chat. Arrange it in the chat. In Telegram chat. Yes. Put it in the
correct order in the Telegram chat. Respond, respond, respond. Have some discussions. Engage. Make
sure you are responding, you are replying to somebody. Agree, disagree. Don't just send your own
response alone. Respond Reply to somebody. Tag somebody. Mention somebody. That's how to discuss.
That's how peer works works. Okay. I'll stop here. All right. The correct answer would be reporter,
sub editor, editor, and then publication. Okay. It's time for us to wrap this up. Check out.

Ask from some question. Ask some questions. Scan it. Scan the QR code. I have question for you. If
you didn't ask me question. And for some of you who have asked asked, I still have question for you.
Where is the highest risk of error? Is it at the What stage is it? At which stage is the most
meaning shift? At which stage is the most Is Is the meaning most shift sort of? What happens if the
sub editor fails in their role? Remember that the sub editor is the key control point. Errors at
this stage multiplies downstream. Yes, if the sub editor didn't catch it, it's likely to, you know,
mean make it into the publication because sometimes the editor just trusts the sub editor.

So, editing is not just correction. It is about quality assurance. Okay. Thank you. Now, it's time
to check out. Check out. Thank you. Check out. All right. If you've been enjoying the lecture, do
well to engage. Like, share, subscribe. Subscribe. Subscribe. And in my last class, we'll be looking
at the news room. What am I saying? We've done this already. Let's check out. I didn't update the
last slide. So, apologies. All right, do well to ask questions where you need to ask questions and
then we'll catch up. Ask your questions there. We'll catch up in the next class or in our live chat.

Cheers and have a good day. Bye-bye.
\end{Verbatim}
}
